% -----------------------------*- LaTeX -*------------------------------
\documentclass[11pt]{report}
\usepackage{../scribe_ds603}

\usepackage{xcolor}

\newcommand{\ab}[1]{\textcolor{purple}{[#1 - AB]}}
\newcommand{\abedit}[1]{\textcolor{purple}{#1}}

\begin{document}




\scribe{Aditya Agrawal}		% required
\lecturenumber{5}			% required, must be a number
\lecturedate{13/08}		% required, omit year


\maketitle


% ----------------------------------------------------------------------
\ab{References to be added throughout}

\section{Input Representations}

We start off by looking at whether the way we represent data for classification
good enough? \ab{this is not just for classification, can mention anomaly
detection as well}
\vspace{1mm}
\\A \emph{feature map} \(\phi :\Rbb^d \ra \Rbb^p\) is simply a map \ab{function}
which maps the input data into another space known as the \emph{feature} or
\emph{embedding} space. 
\vspace{1mm}
\\\abedit{In} $k$-class classification problems, \abedit{we need} a \emph{map}
\(f: X \ra [0,1]^k\) \ab{Make sure the notation for input space is consistent
across scribes} to denote \emph{confidence} of a particular input lying in some
class.
\vspace{1mm}
\\The standard classification algorithms \ab{such as?} directly operate on the input data
while \emph{representational learning} first maps raw data into \emph{feature
space} on which standard algorithms are applied. Representational learning is
generally of two types:

\begin{enumerate}[label=\textup{(\alph*)}, leftmargin=*, widest=b, align=left]



\item \label{mpp:gen_a_0} \emph{Fixed} representations: These are explicit maps
which are used on basis on some prior knowledge of data distribution \ab{give examples!}
\item \label{mpp:gen_a_1} \emph{Learned} representations: These are maps which
are \emph{learnt} through Neural Networks, or other such techniques

 \end{enumerate}
 \vspace{1mm}
A \emph{kernel} is defined as a dot product in feature space:
\begin{align}
    \begin{aligned}
     &k: X \times X \ra \Rbb && k(x,x')=\inprod{\phi(x)}{\phi(x')} \nn
    \end{aligned}   
\end{align}
The standard dot product arises when we use the linear \ab{identity} feature map
\(\phi(x)=x\).
\subsection{Feature maps from kernels}
We ask the question of when do arbitrary kernel functions arise as a dot product
between some feature vectors.
\begin{theorem}
If \(k(x,x')\) is real valued and positive semi-definite, then there exists an
inner product space \(\mathcal{X}\) and feature map \(\phi: \Rbb^d \ra
\mathcal{X}\) such that \(k(x,x')=\inprod{\phi(x)}{\phi(x')}\)
\end{theorem}
\vspace{1mm}
We simply define \(\mathcal{X}\) as the set of all maps \(f: \Rbb^d \ra \Rbb\)
with an appropriate inner product and the feature map \(\phi: \Rbb^d \ra
\mathcal{X}\) as \(\phi(x)(.)=k(x,.)\), which is a function indexed by $x$.\\
\textbf{HW:} Validate the existence of the inner product on the function space
and show that you can recover a PDS kernel.

\subsection{Examples of Kernels}

\begin{enumerate}[label=\textup{(\alph*)}, leftmargin=*, widest=b, align=left]



\item \label{mpp:gen_a_0} Polynomial Kernels : \(k(x,x')= \inprod{x}{x'}^d\)
\item \label{mpp:gen_a_1} Gaussian Kernels :
\(k(x,x')=\epower{-\norm{x-x'}^2/2\sigma^2}\)
 \end{enumerate}


\subsection{The Kernel Trick}
All the representations we dealt with earlier either fixed or learned were into
finite dimensional spaces. But fortunately, in all our optimization equations
\ab{This is only true in the dual formulations}, the mentions of \(\phi(x)\)
were all as dot products of the form \(\phi(x)^T\phi(x')\), which can instead be
replaced by \(\tilde{k}(x,x')\) leading to a feature map \(\tilde{\phi} :
\tilde{\phi}(x)(.)=\tilde{k}(x,.)\) where this feature map could be infinite
dimensional!!

\section{Anomaly Detection via the One-Class SVM}
\ab{Can we get the score notation defined here? It is used in several places below but not explicitly defined}

We come back to the original question of anomaly detection and propose another
method. Consider the following constrained parametric optimization problem:
\ab{Please use bold math for vectors to be consistent with notation in previous scribes}
\begin{align}
        \label{appen:nlmpp}
        \begin{aligned}
            &\minimize_{w,\gamma,\rho}  && \frac{1}{2}\norm{w}^2+\frac{1}{\nu n}(\sum_{i=1}^{n}\gamma_{i})  - \rho\\
            &\sbjto && \begin{cases}
            \inprod{w}{\phi(x_i)} \geq \rho - \gamma_i \quad i=1,\ldots,n,\\
           \gamma_i \geq 0 \quad i=1,\ldots,n,
            \end{cases}
        \end{aligned} \nn
    \end{align}
where \(\nu\) is a hyperparameter which asymptotically represents the fraction
of anomalies. \ab{what are the other terms? How to interpret $\rho$?}
\vspace{1mm}
At test time, we make decisions using \(s,h\) defined as follows:
\begin{align}
        \begin{aligned}
            & s_\theta(x)=\inprod{w}{\phi(x_i)} - \rho\\
            & h_\theta(x)=(-sgn(s_\theta(x))+1)/2
        \end{aligned} \nn
    \end{align}
\textbf{Note:} Finally, the map \(h\) is the one which maps the input into a
class \(0\) or \(1\) where 1 indicates an anomaly.
\noindent \textbf{HW:}
\begin{enumerate}
    \item Derive the primal optimization problem of the \emph{Lagrangian} of a
    standard one class SVM machine for both the soft and the hard cases using
    the KKT conditions
    \item Derive the same for the \(\nu\)-SVM as well
    \item Derive the dual of these primal objective functions as well
\end{enumerate}

\subsection{SLT- like Guarantees on the One Class SVM}
\ab{This is taken from the Learning with kernels book, please cite it}
\begin{theorem}
    (informally) Let \(R_{w,\gamma}=\{x | f_\theta(x) \geq 0\}\). Then for all
    \(\delta > 0 \), with probability \(1- \delta\) over the draws
    \(\{x_i\}_{i=1}^n\), for any \(B >0\), we get
    \begin{align}
        \underset{x \sim \mathbb{P}^+}{\mathbb{P}} [x | x \notin R_{w,\rho-B}] \leq \frac{2}{n}\left( k(n,B,\norm{W},\rho) + \log_2(\frac{n^2}{2\delta}) \right) \nn
    \end{align}
\end{theorem}
where \(k\) is some non-negative function of problem parameters.


\section{Modern OOD Detection}
The OOD problem or the \emph{Out Of Distribution} problem is the supervised case
of the anomaly detection problem \ab{It is not fully supervised, you just have
information about the classes of the in-distribution samples} where we are now
also provided labels on the input draws. Our key question is to understand how
OOD detection is done nowadays. \ab{Please cite the review article on OOD
detection here}

\subsection{Deep OCSVM or Deep SVDD}
\ab{Add citations to the two original papers. They can be found in the Ruff et al. survey paper}
In this methodology, we use a learned feature map \(f_\theta(x)\), where
\(\theta\) parametrizes some neural network, where \(\theta\) is learned jointly
over the supervised learning component (if any) and the anomaly detection \((R,c)\) or
\((w,\gamma,\rho)\). The anomaly detection optimization problem is given as 
\begin{align}
    \begin{aligned}
    &\minimize_{R,c,\theta}  && R^2+\frac{1}{\nu n}\sum_{i=1}^{n}\max (0,\norm{\phi_\theta(x_i)-c}^2-R^2) + \text{Reg}(\theta)
    \end{aligned}  \nn
\end{align}
Here, Reg is some regularization term for \(\theta\) since otherwise it tends to
fit maps such that all points in feature space start falling within the sphere.
It \ab{What is it? Be precise} is typically trained through stochastic gradient
descent.


\subsection{MSP and ODIN}
\ab{Add citations!}
MSP stands for \emph{Maximum Softmax Probability} and works on the principle
that the NN \ab{undefined term. State what functional form the neural network
takes} holds enough discriminative power to classify in distribution samples
confidently. This leads to the formulation where we call an input out of
distribution when the NN cannot classify it confidently. 
\ab{What can you use to check classification confidence?}
\vspace{1mm}
\\ODIN or \emph{Out of DIstribution for Neural networks} enhances MSP by
preprocessing data smartly and adding a temperature control on the softmax
performed. It works in the following way:
\vspace{1mm}
\begin{enumerate}
    \item \label{mpp:gen_0} Preprocesses input \(x\) as
    \begin{align}
        \begin{aligned}
          & x'=x+\eps * sgn \left( \Jacobian_x log(f_\theta^{\hat{y}}(x,\tau)) \right)\\
          & f_\theta^i(x,\tau) = \frac{e^{\frac{l_i(x)}{\tau}}}{\sum_{i=1}^n e^{\frac{l_i(x)}{\tau}}}
        \end{aligned} \nn
    \end{align}
where \(l_i(x)\) is the logit output of neural network \(l: \Rbb^d \ra \Rbb^p \)
\ab{what is $\hat{y}$}


\item \label{mpp:gen_1} Defines the anomaly hypothesis map \(h\) as 
     \begin{align}
        \begin{aligned}
          & s_\theta(x)= \max f_\theta^i(x,\tau) \\
          & h_\theta(x)= \mathbf{1}[s_\theta(x) < \delta] 
        \end{aligned} \nn
    \end{align}
\end{enumerate}

\subsection{Outlier Exposure}
\ab{Cite Deep Anomaly Detection with Outlier Exposure}
The issue with the methods defined so far is that they are overfitted on in
distribution data. It is very likely that out of distribution data looks very
different and might lead the overfit classifier to falsely label them as being
in distribution. To prevent this, we perform outlier exposure by slightly
tweaking our NN loss functions. Our goal is to now minimize
\begin{align}
    \begin{aligned}
            &\minimize_{\theta}  && \underset{x \sim \mathbb{P}^+}{\mathbb{E}} l(h_\theta(x),y) + \lambda  \underset{x \sim \mathbb{P}^-}{\mathbb{E}} l_{OE}(f_\theta(x)) \\
        \end{aligned} \nn
\end{align}
Here \(\lambda\) controls the rate of outliers we expect in testing and
\(\l_{OE}\) is the loss we want to calculate over outlier data. We generally
assume that the out of distribution data is uniform in some sense. \ab{Mention that the OOD data used in OE does not have to match the actual OOD data seen at test time}

We essentially want to temper the confidence of our classifier, by saying "Hey!
Don't be so confident about your predictions, they might be on out of
distribution data!!". We design two outlier loss functions for the same:
\vspace{1mm}
\begin{enumerate}
    \item \label{mpp:gen_0}
    \begin{align}
        l_{OE}(x)= - \sum_{j=1}^{|C|} \frac{1}{|C|} \log(f_\theta^j(x)) \nn
    \end{align}
    We \emph{assume} that a good classification is a confident classification. 
    \ab{here we penalize OOD samples for deviation from the uniform distribution}

    \item \label{mpp:gen_1} To make a distinction between a \emph{good} and
    \emph{confident} classification we artificially introduce a confidence
    estimation branch \(c(x) \in [0,1]\) and set \(l_{OE}(x)=\log(c(x))\).\\
    \vspace{1mm}
    During training we then try to minimize \ab{Cite `Learning Confidence for Out-of-Distribution Detection in Neural Networks'} on the in-distribution data
    \begin{align}
        l(x)= - \sum_{j=1}^{|C|} \log(cf_\theta(x_j)+(1-c)y_j)y_j -B\log(c(x)) \nn
    \end{align}

    Essentially, what we try to say is that if we are not confident about
    whether the input is from distribution, then our NN doesn't pay a training
    loss, however we pay a confidence loss. If we are confident, then we incur a
    training loss as well. This flips our notion from good \(\implies\)
    confident to good \(\impliedby\) confident \ab{Unclear what this means}
    \ab{also mention that the confidence branch provides a score that can be used to classify OOD inputs}
    
\end{enumerate}

\subsection{Grad Norm KL Divergence}
\ab{Cite `On the importance of gradients for detecting distributional shifts in the wild'}
Instead of looking at the maximum softmax probability as a measure of
confidence, we look at the gradient of the norm of the KL divergence \ab{between what?} as such a
measure. The intuition is that low gradient norm implies points close to margins
and so more likely to be OOD. \ab{for inputs that the network is very confident on (implicitly assumed to be in-distribution), you get a large gradient. Why?} We define this using
\begin{align}
    \begin{aligned}
        D_{KL}(U || f_\theta(x)) &= \sum_{j=1}^{|C|} \frac{1}{|C|} \log \left( \frac{1/|C|}{f_\theta^j(x)} \right)\\
        &= - \frac{1}{|C|} \sum_{j=1}^{|C|}  \log \left( |C|f_\theta^j(x) \right)
    \end{aligned} \nn
\end{align}

Our hypothesis function \(h\) then is \(h_\theta(x)= \mathbf{1}
\left[ \norm{\Jacobian_\theta D_{KL}}_2 < \delta \right] \). \ab{What is the score function?}


 

\subsection{Cons of the above methods}
    While the above methods are very intuitive and make sense, most of their
    guarantees are empirical and not theoretical. They have also been designed
    in an adhoc manner with not much theoretical basis as to why they might be
    good. It might be useful to try and get some SLT like bounds on these like
    the one we had for the one class SVM.
% \bibliographystyle{plain} \bibliography{refs}


\end{document}
